\chapter{Technische Erkundung: Proof-of-Concept mit Open Cyber Range (OCR)}
\label{chap:poc-ocr}

Ziel dieser technischen Erkundung war es, neben CyberRangeCZ eine zweite Open-Source-Cyber-Range-Plattform praktisch zu evaluieren.
Im Vordergrund stand dabei kein Produktivbetrieb, sondern ein nachvollziehbarer Proof-of-Concept (PoC), der zeigt, ob sich OCR auf
einem Host deployen lässt und wie die Plattformarchitektur in der Praxis handhabbar ist.

Im Unterschied zu vielen ``All-in-one''-Trainingsplattformen ist OCR nicht als monolithisches Webprodukt konzipiert, sondern als
Baukastensystem aus mehreren Diensten, die gemeinsam den Lebenszyklus einer Übung abdecken: (1) Verwaltung und Distribution von
Übungsartefakten, (2) Orchestrierung und Zustandsverwaltung von Deployments, (3) Anbindung an eine konkrete Infrastruktur
(Virtualisierung/Cloud) sowie (4) Identity- und Access-Management. Diese modulare Architektur ist für Forschung prinzipiell attraktiv,
erhöht im Self-Hosted-Betrieb jedoch den Integrationsaufwand erheblich, da mehrere Komponenten konsistent konfiguriert und miteinander verbunden
werden müssen \cite{OCRDeputyRepo,OCRCZRangerDoc}.

Für den PoC wurden die Kernkomponenten Deputy und die Ranger-Komponente betrachtet. Deputy fungiert als digitale Bibliothek
für Szenarioinhalte: Artefakte (z.\,B.\ Packages, Templates oder vorbereitete Übungskomponenten) werden dort versioniert
abgelegt und können von anderen OCR-Komponenten abgerufen werden.\cite{OCRDeputyRepo,OCRCZFirstDeputyPackage}.

Die Ranger-Komponente bildet die Control Plane für die Ausführung: Sie verwaltet den Zustand einer Übung und koordiniert das
Ausrollen in eine Zielumgebung. Entscheidend ist dabei, dass Ranger die Infrastruktur nicht selbst erzeugt, sondern die Provisionierung
an umgebungsspezifische Handler/Deployer delegiert. OCR bildet damit einen Datenfluss ab: Für OpenStack existiert ein eigener Handler als separates Projekt \cite{OCROpenstackHandler}

\section{Testumgebung}
\label{sec:poc-ocr-testumgebung}

Der PoC wurde auf einer dedizierten Maschine mit Ubuntu 24.04 LTS durchgeführt. Die verwendeten Ressourcen (CPU, RAM und Storage) sind in Abbildung~\ref{fig:poc-ocr-hardware} dargestellt.

\begin{figure}[H]
    \centering
    \includegraphics[width=0.85\linewidth]{poc_ocr_hardware.png}
    \caption{Testumgebung für den OCR-PoC (eigene Aufnahme)}
    \label{fig:poc-ocr-hardware}
\end{figure}

\section{Installation und Deployment der PoC-Umgebung}
\label{sec:poc-ocr-installation}

Die Installation orientierte sich grundsätzlich an der offiziellen Cyber-Range-Dokumentation für die OCR-Komponenten (Deputy und Ranger) \cite{OCRCZRangerDoc,OCRCZFirstDeputyPackage}. In der praktischen Umsetzung kamen aber viele Probleme auf die nicht in der Dokumentation standen: Der Ranger-Service beispielsweiße, prüft beim Start die Konfiguration strikt auf Vollständigkeit und korrekte Datentypen. Fehlende Pflichtfelder (z.\,B.\ für Deployer-/Deployment-Strukturen, ...) führten daher zu wiederholten Startabbrüchen und Neustartschleifen. 

\subsection{Deployment von Deputy}
\label{subsec:poc-ocr-deputy}

Deputy wurde als erster Baustein gestartet, da Ranger für den späteren Betrieb typischerweise Artefakte über Deputy bezieht.
Der Start erfolgte über Docker Compose. Abbildung~\ref{fig:poc-ocr-deputy-up} dokumentiert den erfolgreichen Image-Pull und die laufendenContainer.

\begin{verbatim}
cd /opt/ocr/deputy
docker compose up -d
docker ps
\end{verbatim}

\begin{figure}[H]
    \centering
    \includegraphics[width=1\linewidth]{poc_ocr_deputy_up.png}
    \caption{Deputy: erfolgreicher Start der Services (eigene Aufnahme)}
    \label{fig:poc-ocr-deputy-up}
\end{figure}

\subsection{Deployment der Ranger-Komponente}
\label{subsec:poc-ocr-ranger}

Im nächsten Schritt wurde die Ranger-Komponente deployt. In der Praxis erwies sich dieser Teil als deutlich aufwändiger als erwartet,
da Ranger bereits im Minimalbetrieb strikte Anforderungen an die Konfiguration stellt. Ranger validiert beim Start eine Reihe
verpflichtender Felder und Datentypen. Fehlen diese Angaben oder passen sie nicht exakt, beendet sich der Dienst sofort mit
Fehlermeldungen. 

Der Start und der Containerstatus sind in Abbildung~\ref{fig:poc-ocr-ranger-up} dokumentiert. Abbildung~\ref{fig:poc-ocr-ranger-logs}
zeigt den erfolgreichen Start über Logmeldungen des Webservers (Actix startet Worker).

\begin{figure}[H]
    \centering
    \includegraphics[width=1\linewidth]{poc_ocr_ranger_up.png}
    \caption{Ranger-Komponente(eigene Aufnahme)}
    \label{fig:poc-ocr-ranger-up}
\end{figure}

\begin{figure}[H]
    \centering
    \includegraphics[width=1\linewidth]{poc_ocr_ranger_logs.png}
    \caption{Ranger-Komponente: erfolgreicher Start (eigene Aufnahme)}
    \label{fig:poc-ocr-ranger-logs}
\end{figure}


Im Rahmen des Vorprojekts konnte der Proof-of-Concept für die Open Cyber Range (OCR) in der verfügbaren Zeit nicht bis zu einem lauffähigen Trainings-/Deployment-Start abgeschlossen werden. Obwohl zentrale Komponenten (Deputy und Ranger) grundsätzlich gestartet werden konnten, erwies sich die praktische Integration der modularen OCR-Architektur deutlich komplexer als erwartet, da mehrere Dienste (u.,a. Identity- und Access-Management) konfiguriert und miteinander abgestimmt werden müssen. Zusätzlich setzt die eigentliche Bereitstellung von Übungsumgebungen eine angebundene Virtualisierungs- bzw. Cloud-Infrastruktur voraus (z.,B. eine eigene OpenStack-Umgebung oder alternativ VMware/ESXi über passende Handler/Deployer). Der PoC bleibt daher auf den erfolgreichen Aufbau der Basisservices und die technische Erkundung der Architektur beschränkt.